% 03_advanced_models.tex
% Chapter 3: Advanced Analysis Models [EXTRA]

\section{Advanced Analysis Models}

This chapter defines the optional impairment and diagnostic models added to the clear-sky calculation. Each model can be enabled independently, which also permits the ablation comparison reported in Chapter~5.

\subsection{Adjacent-Satellite Interference (ASI) and IMD}
Finite earth-station beamwidth permits energy from adjacent orbital slots to enter the receive chain and allows an uplink antenna to illuminate neighboring spacecraft. The resulting adjacent-satellite interference (ASI) depends on orbital spacing, antenna pattern, polarization, carrier loading, and relative EIRP. A multicarrier TWTA also generates intermodulation distortion (IMD) when operated close to saturation.

ASI and IMD are included as additional $C/I$ terms so that the clear-sky result is not treated as purely thermal-noise limited.

The off-axis discrimination angle $\theta_{\mathrm{off}}$ between the target satellite and the adjacent interferer is calculated using their respective ECEF position vectors. The off-axis gain $G_{\mathrm{off}}(\theta_{\mathrm{off}})$ in the direction of the interferer is estimated using a simplified internal off-axis antenna-envelope approximation:
\begin{equation}
    G_{\mathrm{off}}(\theta_{\mathrm{off}})
    =
    \min\left[
    G_{\mathrm{max}} - 3.0,\,
    32.0 - 25\log_{10}(\theta_{\mathrm{off}})
    \right]
\end{equation}
This envelope is used for sensitivity analysis, not as an ITU-R regulatory mask or a coordination radiation pattern. The uplink, downlink, and IMD $C/I$ terms are combined with the thermal $C/N$ terms by inverse-linear addition:
\begin{equation}
    \left(\frac{C}{N+I}\right)_{\mathrm{total}}^{-1}
    =
    \left(\frac{C}{N}\right)_{\mathrm{uplink}}^{-1}
    +
    \left(\frac{C}{N}\right)_{\mathrm{downlink}}^{-1}
    +
    \left(\frac{C}{I_{\mathrm{ASI,up}}}\right)^{-1}
    +
    \left(\frac{C}{I_{\mathrm{ASI,down}}}\right)^{-1}
    +
    \left(\frac{C}{I_{\mathrm{IMD}}}\right)^{-1}
\end{equation}

In this project, the ASI/IMD analysis is used as a parametric impairment estimate rather than a regulatory coordination calculation. The adjacent satellites are modeled at $11.0^\circ$E and $15.0^\circ$E, corresponding to a nominal $2.0^\circ$ orbital separation from the target $13.0^\circ$E slot. The adjacent-carrier activity factor is set to $-3.0$~dB, the antenna off-axis rejection follows the simplified envelope above, and the representative transponder IMD term is modeled as $C/I_{\mathrm{IMD}} = 30$~dB. A formal coordination study would require operator-specific carrier plans, measured antenna radiation patterns, polarization reuse information, and transponder loading data.

Chapter~5 reports the resulting margin and separates ASI from the representative IMD term in the ablation table.

\subsection{Hardware Impairments: Non-Linearity and Phase Noise}
The transponder is not an ideal linear amplifier. TWTA compression, AM/PM conversion, and modem imperfections are represented only through simplified engineering terms:
\begin{itemize}
    \item \textbf{Output Back-Off (OBO):} A deliberate reduction of the operating point below saturation to suppress non-linear distortion:
    \begin{equation}
        \mathrm{EIRP}_{\mathrm{active}}
        =
        \mathrm{EIRP}_{\mathrm{sat}}
        -
        \mathrm{OBO}
        \qquad [\mathrm{dBW}]
    \end{equation}
    \item \textbf{Implementation Margin ($M_{\mathrm{impl}}$):} In this study, modem implementation effects are not included as a static RF path loss. Instead, a 1.0~dB DVB-S2 implementation margin is applied only during MODCOD threshold selection:
    \begin{equation}
        \left(\frac{E_s}{N_0}\right)_{\mathrm{usable}}
        =
        \left(\frac{E_s}{N_0}\right)_{\mathrm{available}}
        -
        M_{\mathrm{impl}}
    \end{equation}
\end{itemize}
An OBO value of 2.5~dB is discussed as a representative operating concept \cite{pratt}, but it is not subtracted separately from the published footprint EIRP: the 46~dBW downlink input is treated as the active contour value. Non-linear degradation in the numerical budget is represented by the assumed $C/I_{\mathrm{IMD}}=30$~dB term. A 1.0~dB modem implementation margin is applied only during DVB-S2 MODCOD selection, and the static RF implementation-loss field is set to 0~dB to avoid double counting.

\subsection{Digital Baseband Performance Metrics}
For coherent QPSK over AWGN, the uncoded bit-error probability is:
\begin{equation}
    P_{b,\mathrm{uncoded}}
    =
    \frac{1}{2}\operatorname{erfc}\left(\sqrt{\frac{E_b}{N_0}}\right)
\end{equation}
The report's generic digital-metrics layer applies a 3~dB equivalent coding-gain assumption for sensitivity calculations. The separate ACM calculation does not use this shortcut; it compares the available $E_s/N_0$ directly with the selected DVB-S2 threshold table. Shannon capacity is reported as an upper bound:
\begin{equation}
    C
    =
    B\log_{2}\left(1+\mathrm{SNR}\right)
    \qquad [\mathrm{bps}]
\end{equation}
This coding framework is the foundation for the DVB-S2 MODCOD ladder used by the ACM selector in Chapter~4.

\subsection{Dynamic System Noise Temperature}
The static budget uses a fixed $T_{\mathrm{sys}}$. The dynamic model instead varies the antenna-noise contribution with elevation and with absorptive atmospheric loss. A lossy rain path both attenuates the carrier and increases the apparent sky temperature seen by the receive antenna.

Separating these effects avoids treating rain solely as a carrier-power loss. Chapter~5 reports the carrier attenuation first and then adds the receiver-noise increase as a separate ablation step.

The dynamic, elevation- and fade-dependent system noise temperature $T_{\mathrm{sys}}$ is modeled at the LNA input as:
\begin{equation}
    T_{\mathrm{sys}}(\theta, A_{\mathrm{rain}}) = T_{\mathrm{receiver}} + T_{\mathrm{spillover}} + T_{\mathrm{sky}}(\theta) + T_{\mathrm{rain}}(A_{\mathrm{rain}})
\end{equation}
where $T_{\mathrm{sky}}(\theta)$ is the clear-sky noise at elevation angle $\theta$, and $T_{\mathrm{rain}}$ is the rain thermal emission noise calculated using the radiative transfer equation:
\begin{equation}
    T_{\mathrm{rain}} = T_{\mathrm{atm}} \left( 1 - 10^{-A_{\mathrm{rain}}/10} \right)
\end{equation}
where $T_{\mathrm{atm}} = 275$~Kelvin represents the physical atmospheric temperature of the lossy rain medium.

For example, 10~dB of absorptive loss gives $T_{\mathrm{rain}}\approx247.5$~K for $T_{\mathrm{atm}}=275$~K. This is the rain-emission contribution only; the total $T_{\mathrm{sys}}$ also depends on the receiver, spillover, and clear-sky terms and is calculated from the full expression above.

\subsection{Apparent GEO Orbit Motion}
The apparent spacecraft position is allowed to vary within a simplified station-keeping box:
\begin{align}
    \Delta \lambda(t) &= A_{\mathrm{ew}} \sin\left(\frac{2\pi t}{T_{\mathrm{sidereal}}} + \phi_{\mathrm{ew}}\right) \\
    \Delta \phi(t) &= A_{\mathrm{ns}} \sin\left(\frac{2\pi t}{T_{\mathrm{sidereal}}} + \phi_{\mathrm{ns}}\right) \\
    \Delta r(t) &= A_{\mathrm{r}} \sin\left(\frac{2\pi t}{T_{\mathrm{sidereal}}} + \phi_{\mathrm{r}}\right)
\end{align}
These offsets update the slant-range vector $\vec{\rho}(t)$. The model is intended to test sensitivity to a controlled GEO station-keeping box; it does not establish a tracking requirement, which would depend on the measured antenna pattern, pointing tolerance, and operator practice.
